% !TEX root = ../main.tex
\chapter{Ausblick (Future Work)}
Der aktuelle Prototyp des Assistenzsystems \glqq Ascenta\grqq\ belegt die Machbarkeit einer dezentralen, offline-fähigen Smart-Glasses-Architektur. Um das System jedoch von einem funktionalen Prototyp zu einem marktreifen Endkundenprodukt weiterzuentwickeln, sind in zukünftigen Iterationen verschiedene Erweiterungen und Optimierungen in allen Fachbereichen erforderlich.

\section{Software und Benutzeroberfläche - Liam}
Die zukünftige Software-Entwicklung muss sich primär auf die vollständige Barrierefreiheit und erweiterte Navigationskonzepte konzentrieren.

\subsection{Native Zero-UI-Applikation für blinde Nutzer}
Die aktuelle Android-Applikation verfügt über eine grafische Benutzeroberfläche (GUI), die essenziell für das Debugging, die Evaluierung und die Ersteinrichtung durch sehende Begleitpersonen ist. In einer finalen Version für den Endanwender muss diese App jedoch als reine Hintergrundanwendung (\textit{Zero-UI}) konzipiert werden. 

Das visuelle Interface wird dabei vollständig obsolet. Die Applikation sollte sich nahtlos in native Accessibility-Dienste (wie Android TalkBack) integrieren und den Setup-Prozess (z.\,B. die Eingabe des WLAN-Passworts) über auditive Menüführung oder NFC-Tags abwickeln. Der Fokus liegt dabei auf maximaler Stabilität des Foreground-Services und einer Reduktion des Batterie- und Speicherverbrauchs des Smartphones, da keine visuellen Elemente gerendert werden müssen.

\subsection{Kontinuierliches auditives Raum-Feedback (Leitstrahl-Prinzip)}
Eine vielversprechende funktionale Erweiterung stellt die Implementierung eines dynamischen, akustischen Suchmodus dar. Sucht der Nutzer nach einem bestimmten Objekt (z.\,B. einem Stuhl oder einer Tür), könnte das System ein kontinuierliches Audiosignal generieren, das an die Funktion eines Geigerzählers oder eines Einparkhilfesensors angelehnt ist. 



Wenn der Nutzer den Kopf schwenkt und das gesuchte Objekt in das Sichtfeld der Kamera gerät, passt der Algorithmus die Frequenz (Tonhöhe) und das Intervall (Pings pro Sekunde) des Audiosignals in Echtzeit an. Je zentrierter das Objekt im Sichtfeld ist und je geringer die Distanz (gemessen durch den Time-of-Flight-Sensor) wird, desto intensiver wird das akustische Feedback. Dies würde blinden Nutzern eine hochpräzise, intuitive Zielführung ermöglichen, ohne sie mit komplexen, gesprochenen Sätzen zu überlasten.

\section{Künstliche Intelligenz und Modell-Training}

\subsection{Erweiterte Objekterkennung für Alltagsgegenstände}
\label{sec:future_ml_objects}
Derzeit ist das System auf die Erkennung von Personen optimiert. Für eine umfassende Alltagshilfe muss das zugrundeliegende Dataset massiv erweitert werden. Die Firmware und die App-Logik sind bereits so abstrahiert, dass sie beliebige Labels verarbeiten können. Zukünftige Iterationen müssen daher trainiert werden, um eine Vielzahl von Hindernissen und Zielen (z.\,B. Treppenstufen, Ampeln, Zebrastreifen, Türen, Sitzgelegenheiten oder herannahende Fahrzeuge) zuverlässig zu klassifizieren.

\subsection{Weiterentwicklung des Machine-Learning-Modells (Edge Impulse)}
\label{sec:future_edge_impulse}
Fabian

\section{Hardware und Elektronik}

\subsection{Optimierung der Schaltungstechnik und Sensorik}
\label{sec:future_elektronik}
Fabian

\section{Mechanik und Gehäusedesign}

\subsection{Weiterentwicklung des 3D-Modells und der Ergonomie}
\label{sec:future_3ddruck}
Manuel